\section*{คำชี้แจงสำหรับผู้เข้าร่วม}

โปรดตอบตามประสบการณ์จากการทดลองใช้รอบนี้ ทั้งข้อดีและปัญหาที่พบ คำตอบใช้ปรับปรุงแอป ไม่มีคำตอบถูกหรือผิด ไม่ต้องระบุชื่อจริง ไม่ต้องระบุข้อมูลสุขภาพส่วนตัวในช่องข้อเสนอแนะ การเข้าร่วมเป็นความสมัครใจและหยุดได้ ก่อนเริ่มผู้ดำเนินการต้องชี้แจงผู้รับผิดชอบ ช่องทางติดต่อ ข้อมูลที่เก็บ ผู้เข้าถึงข้อมูล ระยะเวลาเก็บรักษา และการขอถอนข้อมูลตามการดำเนินงานจริง

\noindent\textbf{การยินยอม}\quad
$\square$ ได้รับคำชี้แจงและยินยอมเข้าร่วม\quad
$\square$ ไม่ยินยอม (จบแบบสอบถาม)

\medskip
\noindent ระบบออนไลน์บันทึกวันที่และเวลาส่งคำตอบอัตโนมัติ

\noindent อายุ 18 ปีขึ้นไป\quad $\square$ ใช่\quad $\square$ ไม่ใช่

\noindent ก่อนหน้านี้ ท่านเคยใช้ Pacegasus หรือไม่\quad $\square$ ไม่เคย\quad $\square$ เคย

\noindent ระยะเวลาวิ่งต่อเนื่องโดยประมาณ \dotfill\ เดือน

\noindent ฟีเจอร์ที่ได้ลองใช้ในรอบนี้ (เลือกได้หลายข้อ)

\noindent $\square$ เป้าหมายและตารางฝึก\quad $\square$ Wellness Check-in

\noindent $\square$ บันทึกกิจกรรม\quad $\square$ ผลสรุปและ RPE\quad $\square$ ภารกิจและรางวัล

\noindent $\square$ คำแนะนำการฝึก\quad $\square$ เพื่อน คลับ หรือการจัดอันดับ\quad $\square$ อื่น ๆ \dotfill

\section*{ส่วนที่ 1 ความเข้าใจและประสบการณ์การใช้งาน}

โปรดเลือกหนึ่งคำตอบต่อข้อ โดย 1 = ไม่เห็นด้วยอย่างยิ่ง, 2 = ไม่เห็นด้วย, 3 = ไม่แน่ใจ, 4 = เห็นด้วย และ 5 = เห็นด้วยอย่างยิ่ง ช่อง ``N/A'' ในแบบกระดาษหมายถึงไม่ได้ทดลอง/ประเมินไม่ได้ ส่วนฉบับออนไลน์ให้เว้นข้อแทน หากได้ลองแล้วแต่ทำไม่สำเร็จ สามารถให้คะแนนตามประสบการณ์นั้นได้

\begingroup
\small
\setlength{\tabcolsep}{3pt}
\renewcommand{\arraystretch}{1.25}
\begin{longtable}{|>{\raggedright\arraybackslash}p{0.085\textwidth}|>{\raggedright\arraybackslash}p{0.525\textwidth}|*{5}{>{\centering\arraybackslash}p{0.035\textwidth}|}>{\centering\arraybackslash}p{0.060\textwidth}|}
\hline
\textbf{รหัส} & \textbf{ข้อความประเมิน} & \textbf{1} & \textbf{2} & \textbf{3} & \textbf{4} & \textbf{5} & \textbf{N/A} \\ \hline
\endfirsthead
\hline
\textbf{รหัส} & \textbf{ข้อความประเมิน} & \textbf{1} & \textbf{2} & \textbf{3} & \textbf{4} & \textbf{5} & \textbf{N/A} \\ \hline
\endhead
U01 & ฉันเข้าใจวิธีใช้งานเบื้องต้นของแอปได้ง่าย & $\square$ & $\square$ & $\square$ & $\square$ & $\square$ & $\square$ \\ \hline
U02 & ฉันค้นหาฟังก์ชันที่ต้องการในแอปได้ง่าย & $\square$ & $\square$ & $\square$ & $\square$ & $\square$ & $\square$ \\ \hline
U03 & ขั้นตอนเลือกเป้าหมายการวิ่งเข้าใจได้ง่าย & $\square$ & $\square$ & $\square$ & $\square$ & $\square$ & $\square$ \\ \hline
U04 & ฉันจัดวันฝึกตามเวลาที่ตนสะดวกได้ง่าย & $\square$ & $\square$ & $\square$ & $\square$ & $\square$ & $\square$ \\ \hline
U05 & ฉันเข้าใจคำถามใน Wellness Check-in & $\square$ & $\square$ & $\square$ & $\square$ & $\square$ & $\square$ \\ \hline
U06 & ฉันเริ่มบันทึกกิจกรรมได้สะดวก & $\square$ & $\square$ & $\square$ & $\square$ & $\square$ & $\square$ \\ \hline
U07 & ฉันจบและบันทึกกิจกรรมได้สะดวก & $\square$ & $\square$ & $\square$ & $\square$ & $\square$ & $\square$ \\ \hline
U08 & ข้อมูลสรุปหลังจบกิจกรรมเข้าใจได้ง่าย & $\square$ & $\square$ & $\square$ & $\square$ & $\square$ & $\square$ \\ \hline
U09 & คำอธิบายระดับความเหนื่อยช่วยให้ฉันเลือกคะแนน RPE ได้ & $\square$ & $\square$ & $\square$ & $\square$ & $\square$ & $\square$ \\ \hline
U10 & ฉันเข้าใจว่ารางวัลที่แสดงได้รับจากกิจกรรมใด & $\square$ & $\square$ & $\square$ & $\square$ & $\square$ & $\square$ \\ \hline
E01 & ภารกิจที่ได้ทดลองทำให้การใช้แอปน่าสนใจสำหรับฉัน & $\square$ & $\square$ & $\square$ & $\square$ & $\square$ & $\square$ \\ \hline
E02 & การแสดงผลช่วยให้ฉันเห็นความก้าวหน้าของกิจกรรมในรอบนี้ & $\square$ & $\square$ & $\square$ & $\square$ & $\square$ & $\square$ \\ \hline
\end{longtable}
\endgroup

\noindent\begin{minipage}{\linewidth}
\section*{ส่วนที่ 2 ความพึงพอใจและความตั้งใจใช้ต่อ}

\noindent\textbf{S01} ท่านพึงพอใจกับการใช้ Pacegasus ในรอบนี้มากเพียงใด

\noindent $\square$ 1 น้อยที่สุด\quad $\square$ 2 น้อย\quad $\square$ 3 ปานกลาง

\noindent $\square$ 4 มาก\quad $\square$ 5 มากที่สุด\quad $\square$ ประเมินไม่ได้

\medskip
\noindent\textbf{I01} หลังทดลองครั้งนี้ ท่านตั้งใจจะใช้ Pacegasus ต่อมากเพียงใด

\noindent $\square$ 1 น้อยที่สุด\quad $\square$ 2 น้อย\quad $\square$ 3 ปานกลาง

\noindent $\square$ 4 มาก\quad $\square$ 5 มากที่สุด\quad $\square$ ยังตัดสินใจไม่ได้
\end{minipage}\par

\section*{ส่วนที่ 3 ฟีเจอร์ที่ได้ทดลองเพิ่มเติม}

ส่วนนี้ตอบเฉพาะผู้ที่ใช้ฟีเจอร์นั้นจริงในรอบนี้ ใช้สเกลเห็นด้วย 1--5 เช่นเดียวกับส่วนที่ 1 และ \mbox{N/A} สำหรับประเมินไม่ได้ หากไม่เคยเห็นคำแนะนำหรือไม่เคยใช้เพื่อน คลับ หรือการจัดอันดับ ให้ข้ามข้อที่เกี่ยวข้อง

\begingroup
\small
\setlength{\tabcolsep}{3pt}
\renewcommand{\arraystretch}{1.25}
\begin{longtable}{|>{\raggedright\arraybackslash}p{0.085\textwidth}|>{\raggedright\arraybackslash}p{0.525\textwidth}|*{5}{>{\centering\arraybackslash}p{0.035\textwidth}|}>{\centering\arraybackslash}p{0.060\textwidth}|}
\hline
\textbf{รหัส} & \textbf{ข้อความประเมิน} & \textbf{1} & \textbf{2} & \textbf{3} & \textbf{4} & \textbf{5} & \textbf{N/A} \\ \hline
O01 & คำแนะนำการฝึกที่ได้รับในรอบนี้ดูเหมาะสมกับความพร้อมของฉัน & $\square$ & $\square$ & $\square$ & $\square$ & $\square$ & $\square$ \\ \hline
O02 & จากการทดลองรอบนี้ ฉันไว้วางใจคำแนะนำการฝึกที่ระบบแสดง & $\square$ & $\square$ & $\square$ & $\square$ & $\square$ & $\square$ \\ \hline
O03 & ฟีเจอร์เพื่อน คลับ และการจัดอันดับที่ได้ลอง ทำให้ฉันรู้สึกได้รับการสนับสนุนในการวิ่ง & $\square$ & $\square$ & $\square$ & $\square$ & $\square$ & $\square$ \\ \hline
\end{longtable}
\endgroup

\section*{ส่วนที่ 4 ข้อเสนอแนะจากการทดลองรอบนี้}

\noindent\textbf{F01} พบจุดใดที่สับสนหรือทำต่อไม่ได้ เกิดที่หน้าจอใด

\noindent\dotfill\\[0.5\baselineskip]\noindent\dotfill

\medskip
\noindent\textbf{F02} ฟีเจอร์ใดมีประโยชน์หรือน่าสนใจที่สุด เพราะเหตุใด

\noindent\dotfill\\[0.5\baselineskip]\noindent\dotfill

\medskip
\noindent\textbf{F03} หากปรับปรุงได้หนึ่งอย่าง ท่านอยากให้แก้อะไร

\noindent\dotfill\\[0.5\baselineskip]\noindent\dotfill

\clearpage
\section*{แบบบันทึกภารกิจสำหรับผู้ดำเนินการ}

\noindent รหัสผู้เข้าร่วม \dotfill\quad วันที่ \dotfill

\noindent รุ่นแอป \dotfill\quad รุ่นโทรศัพท์/Android \dotfill

\noindent เวลาเริ่มรอบ \dotfill\quad เวลาจบรอบ \dotfill

\noindent โหมดทดลอง\quad $\square$ ใช้งานจริง\quad $\square$ จำลอง\quad $\square$ มีทั้งสองแบบ

อ่านสถานการณ์จากตารางภารกิจในบทที่ 3 หรือคู่มือสร้างฟอร์มที่แนบกับชุดเครื่องมือ ให้ผู้ใช้ลองเองก่อน ไม่สอนตำแหน่งปุ่มล่วงหน้า บันทึกการช่วยทุกครั้ง หากต้องข้ามภารกิจ ให้ระบุสาเหตุและเตรียมสถานะเริ่มต้นของภารกิจถัดไปโดยไม่เปลี่ยนผลของภารกิจก่อนหน้า

\noindent สถานะ: A = สำเร็จเอง, H = สำเร็จเมื่อช่วย, F = ไม่สำเร็จ, N = ไม่ได้ทดลอง, D = สาธิตอย่างเดียว

\begingroup
\small
\setlength{\tabcolsep}{4pt}
\renewcommand{\arraystretch}{1.4}
\begin{longtable}{|>{\raggedright\arraybackslash}p{0.07\textwidth}|>{\raggedright\arraybackslash}p{0.29\textwidth}|>{\centering\arraybackslash}p{0.11\textwidth}|>{\centering\arraybackslash}p{0.11\textwidth}|>{\raggedright\arraybackslash}p{0.25\textwidth}|}
\hline
\textbf{รหัส} & \textbf{ภารกิจ} & \textbf{สถานะ} & \textbf{เวลา (วินาที)} & \textbf{ปัญหา/การช่วย} \\ \hline
T01 & เข้าระบบและเตรียมข้อมูลเริ่มต้น & & & \\[1.3\baselineskip] \hline
T02 & เป้าหมายและวันฝึก & & & \\[1.3\baselineskip] \hline
T03 & บันทึกความพร้อม & & & \\[1.3\baselineskip] \hline
T04 & เริ่มและจบกิจกรรม & & & \\[1.3\baselineskip] \hline
T05 & RPE และอ่านผลสรุป & & & \\[1.3\baselineskip] \hline
T06 & ความก้าวหน้าและรางวัล & & & \\[1.3\baselineskip] \hline
\end{longtable}
\endgroup

เวลานับจากผู้เข้าร่วมเริ่มลงมือหลังเข้าใจโจทย์จนถึงสำเร็จหรือหยุดภารกิจ ให้จดเวลาเริ่ม--สิ้นสุด จำนวนครั้งที่ช่วย และเหตุขัดข้องเพิ่มเติม หากมีเวลารอ OTP หรือเครือข่าย ให้ระบุแยกเพื่อช่วยตีความ ผู้ที่ทำไม่สำเร็จยังคงอยู่ในจำนวนผู้ที่ได้ลองภารกิจ การสาธิตอย่างเดียวไม่นับเป็นผู้เข้าร่วมทำสำเร็จ

\section*{แนวทางสนทนาหลังทดลองใช้}

ถามต่อจากคำตอบ F01--F03 ว่าเกิดอะไรขึ้น ผู้ใช้คาดว่าจะพบอะไร และผลกระทบเป็นอย่างไร หากเกี่ยวกับคำแนะนำให้ถามว่าเหตุใดจึงรู้สึกเหมาะสมหรือไม่เหมาะสม โดยไม่ชี้นำว่าระบบช่วยให้วิ่งดีขึ้นแล้ว บันทึกข้อสังเกตเชื่อมกับรหัสภารกิจและรหัสผู้เข้าร่วม

\clearpage
\section*{ที่มาของประเด็นและวิธีสรุปผล}

\begingroup
\small
\renewcommand{\arraystretch}{1.15}
\begin{longtable}{|>{\raggedright\arraybackslash}p{0.19\textwidth}|>{\raggedright\arraybackslash}p{0.36\textwidth}|>{\raggedright\arraybackslash}p{0.35\textwidth}|}
\hline
\textbf{เครื่องมือ} & \textbf{ที่มาของกรอบ} & \textbf{ขอบเขตการตีความ} \\ \hline
\endfirsthead
\hline \textbf{เครื่องมือ} & \textbf{ที่มาของกรอบ} & \textbf{ขอบเขตการตีความ} \\ \hline
\endhead
T01--T06 & การทำภารกิจและสังเกตตาม Moran; สถานการณ์ตาม McCloskey \cite{moran2019usability,mccloskey2014tasks} & ผลสำเร็จ ปัญหา และการช่วยในรอบนี้ \\ \hline
เวลา/สถานะ & แนวทางเก็บผลสำเร็จ เวลา และข้อผิดพลาดของ NIST \cite{nistUsabilityTesting} & สรุปรายคนและรายภารกิจ ไม่ยืนยันค่าของประชากร \\ \hline
U01--U10 & การเรียนรู้ครั้งแรกและความสะดวกตาม Nielsen \cite{nielsen2012usability}; เส้นทางฟังก์ชันจริงของ Pacegasus & ความเข้าใจที่ผู้ใช้รายงาน ไม่ใช่การตรวจความแม่นยำของระบบ \\ \hline
E01--E02 & Gamification และการรับรู้ความก้าวหน้า เชื่อมกรอบ SDT ในบทที่ 2 \cite{ryan2020} & ความน่าสนใจและความก้าวหน้าที่รับรู้ครั้งแรก ไม่ใช่แบบวัดแรงจูงใจมาตรฐาน \\ \hline
S01 และ I01 & ความพึงพอใจตาม NIST \cite{nistUsabilityTesting}; ความตั้งใจใช้ต่อเป็นข้อเพิ่มเติมของโครงงาน & รายงานแยก ความตั้งใจไม่ใช่ retention \\ \hline
O01--O03 & ประเด็นออกแบบคำแนะนำและระบบสังคมของโครงงาน; SDT สำหรับความสัมพันธ์ \cite{ryan2020} & เฉพาะผู้ที่ได้ลองจริง ไม่ยืนยันผลด้านสุขภาพหรือความปลอดภัย \\ \hline
F01--F03 & การรับฟังข้อเสนอแนะควบคู่การสังเกต \cite{moran2019usability} & ปัญหา เหตุผล และข้อเสนอเพื่อปรับปรุง \\ \hline
\end{longtable}
\endgroup

แบบประเมินใช้เก็บประสบการณ์ทันทีหลังทดลองหนึ่งรอบ การอ้างอิงในตารางรองรับกรอบและวิธีเก็บข้อมูล ส่วนข้อความภาษาไทย สเกล และรหัสข้อเป็นสิ่งที่ผู้จัดทำออกแบบ ยังไม่ใช่เครื่องมือที่ผ่านการตรวจสอบความตรงหรือความเชื่อมั่น ขนาดกลุ่มประมาณ 5 คนใช้ค้นหาปัญหาเบื้องต้น ไม่ใช่เพื่อยืนยันคะแนนแทนประชากร \cite{budiu2021five}

สรุปคะแนนแยกรายข้อ โดยแสดงจำนวนคำตอบที่ใช้ได้ ค่ามัธยฐาน ช่วงต่ำสุด--สูงสุด และอาจเพิ่มค่าเฉลี่ยเชิงพรรณนา แยก N/A และยังตัดสินใจไม่ได้ออกจากคะแนน ไม่สร้างคะแนนรวม SUS หรือคะแนนประสิทธิผลระยะยาว สรุปผลภารกิจจาก A, H และ F พร้อมจำนวนผู้ได้ลองจริง; แสดง N และ D แยกต่างหาก

\section*{การนำไปใช้ใน Google Forms และ Google Sheets}

สร้างฟอร์มหนึ่งชุด ใช้สเกลเชิงเส้นสำหรับคะแนน 1--5 พร้อมคำอธิบายแต่ละระดับ ปิดการบังคับตอบข้อคะแนนเพื่อให้เว้นได้เมื่อประเมินไม่ได้ และใช้ย่อหน้าสำหรับข้อเสนอแนะ ไม่เก็บรหัสและวันที่จากผู้ตอบ โดยใช้เวลาส่งคำตอบอัตโนมัติ ใช้คำถามทางแยกเพื่อแสดง O01--O03 เฉพาะฟีเจอร์ที่ได้ลองจริง หากไม่ยินยอมให้จบการตอบ และไม่จำกัดหนึ่งคำตอบต่อคอลัมน์ของตารางคะแนน \cite{googleFormsQuestionTypes}

เชื่อมคำตอบไปยัง Google Sheets \cite{googleFormsResponses} เก็บข้อมูลดิบแยกจากผลรายภารกิจ สรุปรายข้อ และประเด็นปัญหา บันทึกรุ่นแอป โหมดทดลอง สถานะ เวลา และการช่วย ส่วนผลรายข้อแสดงจำนวนคำตอบที่ใช้ได้ \mbox{N/A} และข้อมูลขาดแยกกัน
